% SHARD-ID: RC-04U-DENINGER-LPS
% SHARD-TITLE: The Phantasm XI: the (13, 17; 5) LPS Square Complex as a Deninger Finite Model. Horizontal Cohomology is the H^1 of the 13-Graph, the Transverse Transport is T_17, the Spectrum is a Weight-Two Hecke Spectrum, the Bass Doubling is Frobenius at 17 as a One-Operator System, the Metric Cone Has Dimension 16539, and the Weil--LPS Channel Is a Different Selection of the Same Hecke Family
% SHARD-SUMMARY: The square complex of the worked example is Gamma(5) \ (T_14 x T_18) for the primary quaternion product lattice (Mozes, Rattaggi), with total Betti numbers (1, 0, 5759) certified by an exact rank; its 721-dimensional horizontal cohomology is H^1 of the LPS graph Lambda_13(5) \ T_14 and the square transport is the degree-18 T_17 correspondence, so chi_T = chi_{T_1}(x - 18)/chi_{A_17}; the directed new space is K (+) K' with the bipartite sign intertwining T_17 with -T_17 on the other half.
% SHARD-SUMMARY: The vertex functions are algebraic automorphic forms on the Hamilton quaternion algebra of level U_2 = image(1 + 2 O_2) at 2 and principal level at 5; K is the Steinberg-at-13 newspace with reversal sign -1, T = a_17, multiplicities dim pi_2^{U_2} dim pi_5^{U_5}, and |a_17| <= 2 sqrt 17 is Eichler--Shimura--Igusa; as a PGL_2(F_5)-module K = 6 Reg + 1. The Bass doubling (K (+) K, F_1) is the direct sum of the Frobenius-at-17 two-dimensional Galois systems as one-operator systems; a canonical cup/Petersson comparison is not supplied by this, the certificate G differs from the compatible metric Omega_B J = G/r(T), and the real cone of invariant positive metrics has dimension sum m_a^2 = 16539.
% SHARD-SUMMARY: The notebook's (13; 17) Weil--LPS channel is the Fourier-block selection of the same spherical Hecke family at principal level 13 (End(W_-) = 1 + 7 + 14 + 14 contains no Steinberg type); its exact odd characteristic polynomial shares only (x - 3)^8 (x + 3)^3 (x + 5)^6 with chi_T and neither divides the other. An exact graded CP realisation of the net divisor is open: the direct sector prescription has Tr E^2 = -10940, and scalar signs cancel under Ad; a unitary-memory CP Hashimoto transfer exists. The two-prime orbit quotient is non-Hausdorff.
% SHARD-KEYWORDS: LPS graphs, square complex, product of trees, quaternion lattice, Mozes, Rattaggi, horizontal cohomology, Hecke correspondence, Jacquet--Langlands, Steinberg, Iwahori, Eichler--Shimura, Bass doubling, companion matrix, invariant metric cone, Weil representation, Gerardin, graded CP transfer, Kraus, Hashimoto
% SHARD-NOTES: none
% SHARD-SCRIPTS: scripts/lps_square_complex.py

\section{The Phantasm XI: the LPS Square Complex as a Deninger Finite Model}
\label{sec:deninger-lps}

Round of 2026-09-22 (brief \texttt{notes/deninger-lps/astra-brief.md}, prover \texttt{notes/deninger-lps/astra-proofs.md},
blind numerics \texttt{scripts/lps\_square\_complex.py}, review \texttt{notes/reviews/deninger-lps-2026-09-22.md}). The
input is a worked example from a separate Codex conversation on Deninger's programme (\texttt{notes/deninger-lps/src/}):
vertex set $\mathrm{PGL}_2(\mathbb F_5)$ ($120$), horizontal generators $S_{13}$ ($14$), vertical generators $S_{17}$ ($18$),
$840+1080$ edges, $7560$ squares from the $252$ reordering rules $ab = \pm b'a'$; horizontal differential $D$, horizontal
cohomology $K = \ker D^{t}$ of dimension $721$, transverse transport $T_1$ through the squares with $T_1D = DA_{17}$, $T = T_1|_K$
with $\chi_T$ fully factored ($\|T\| = 4+2\sqrt3<8<2\sqrt{17}$); Bass doubling $F_1 = \begin{psmallmatrix}T&-17\\1&0\end{psmallmatrix}$ on
$K\oplus K$, pairing $\Omega_B = \begin{psmallmatrix}0&1\\-1&0\end{psmallmatrix}$, certificate $G = \begin{psmallmatrix}1&-T/2\\-T/2&17\end{psmallmatrix}>0$, % bare-ok
quarter-turn $J$, graded Euler product with the $(1-u^2)^{5760}$ correction; the example's author left the ``face-derived
identification'' of the pairing open. Here $\Omega$ denotes an alternating pairing. % bare-ok

\subsection{The complex and its horizontal cohomology}

\begin{proposition}[The complex is a congruence cover of the primary quaternion product lattice]
\label{prop:lps-square-complex-structure}
\claimstatus{prop:lps-square-complex-structure}{proved}
Let $B = (-1,-1)_{\mathbb Q}$ (a division algebra, not $M_2(\mathbb Q)$) and $\Gamma = \langle[S_{13}],[S_{17}]\rangle\subset PB^\times(\mathbb Q)$,
embedded diagonally in $PB^\times(\mathbb Q_{13})\times PB^\times(\mathbb Q_{17})$; by Mozes (1995) and Rattaggi (2004), not
byte-cited, the projective group generated by the primary integral quaternions of norms $13$ and $17$ acts freely and
transitively on the vertices of $T_{14}\times T_{18}$ (the full projective unit group has torsion and is not this lattice).
With $\Gamma(5) = \ker(\Gamma\to\mathrm{PGL}_2(\mathbb F_5))$ (onto: the horizontal generators alone generate the $120$ elements;
nonsquare norms alone prove only the determinant parity), $X = \Gamma(5)\backslash(T_{14}\times T_{18})$ has vertices
$\Gamma(5)\backslash\Gamma$ and exactly the edges and squares of the worked example. Total cohomology: $\mathrm{rank}\,B_1 = 119$,
$\mathrm{rank}_{\mathbb F_{101}}(B_2B_2^{t}) = 1801$ and $B_1B_2 = 0$ give $b_1 = 0$, $b_2 = 5759$ exactly. No general vanishing theorem
is used: $F_r\times F_s$ on $T_{2r}\times T_{2s}$ has $b_1\ne0$.
\end{proposition}

\begin{theorem}[Horizontal cohomology is graph $H^1$; the square transport is the $T_{17}$ correspondence]
\label{thm:lps-horizontal-h1-hecke}
\claimstatus{thm:lps-horizontal-h1-hecke}{proved}
Fix a vertex $o_{17}$ of $T_{18}$, $\Lambda_{13} = \mathrm{Stab}_\Gamma(o_{17}) = \langle S_{13}\rangle$, $\Lambda_{13}(5) = \Lambda_{13}\cap\Gamma(5)$.
The horizontal colour graph is $Y = \Lambda_{13}(5)\backslash T_{14}$ (not $\Gamma(5)\backslash T_{14}$: $\Gamma(5)$ is a lattice in the
product), connected, so $K = \ker\partial = H_1(Y;\mathbb R)\cong H^1(Y;\mathbb R)$, $\dim K = 840-120+1 = 721$. The correspondence
$Y\xleftarrow{p}C\xrightarrow{r}Y$ with $C$ the graph of directed vertical edges and squares gives $p_{*}r^{*}f(g,a) =
\sum_{b\in S_{17}}f(g\phi(b'),a')$, $ab = \pm b'a'$: this is $T_1$, self-adjoint (reversing vertical edges exchanges $p,r$), commuting
with horizontal reversal $R$ on directed edge fields (the antisymmetric convention is the $R = -1$ sector), the spherical
double coset at $17$ of degree $18$, i.e. the unnormalised $T_{17}$; telescoping around a square gives $T_1D = DA_{17}$, the
$252$ rules give $A_{13}A_{17} = A_{17}A_{13}$ at edge level, and $C^1 = \operatorname{im}D\oplus K$ gives
$\chi_T = \chi_{T_1}(x)(x-18)/\chi_{A_{17}}(x)$. On directed edges the old space $\operatorname{im}(s^{*},t^{*})$ has dimension
$240-2 = 238$ (kernel spanned by $(1,-1)$ and $(\varepsilon,\varepsilon)$, $\varepsilon$ the bipartite sign), the new space $N$ has
$1442 = 2\cdot721$, $N_- = K$, and $Sf(g,a) = \varepsilon(g)f(g,a)$ satisfies $SR = -RS$, $ST_1 = -T_1S$, $S: N_-\to N_+$: the other
half carries $-\mathrm{spec}\,T$ (traces $18$ and $-18$), $\chi_{T_1|N}(x) = \chi_T(x)\cdot(-1)^{721}\chi_T(-x)$. Verified:
\cref{num:lps-square-complex}.
\end{theorem}

\begin{theorem}[The spectrum of $T$ is a weight-two Hecke spectrum at $17$ (Jacquet--Langlands, Casselman, Eichler--Shimura--Igusa)]
\label{thm:lps-spectrum-hecke}
\claimstatus{thm:lps-spectrum-hecke}{proved}
With $O$ the Hurwitz order (norm-Euclidean, class number one; projective units of order $12$ identified with $(O/2O)^\times$),
$U_2 = $ image of $1+2O_2$ in $PB^\times(\mathbb Q_2)$, $U_5 = \ker(\mathrm{PGL}_2(\mathbb Z_5)\to\mathrm{PGL}_2(\mathbb F_5))$, $U_l = \mathrm{PGL}_2(\mathbb Z_l)$
otherwise: the vertex functions are $\mathcal A(U) = \mathbb C[G(\mathbb Q)\backslash G(\mathbb A_f)/U]$, weight zero on the compact real
group (holomorphic weight two after Jacquet--Langlands); ``maximal at $2$'' would quotient by nontrivial Hurwitz units.
Replacing $U_{13}$ by the ordered-edge stabiliser gives $\mathbb C(E^\to)$; by Casselman (1973) an irreducible generic
representation of $\mathrm{PGL}_2(\mathbb Q_{13})$ with Iwahori invariants is spherical (two-dimensional invariants, old) or
$\mathrm{St}\otimes\chi$, $\chi$ unramified quadratic (one-dimensional, reversal eigenvalue $-\chi(13)$). Hence, over $\mathbb C$ with all
coefficient embeddings, $K_{\mathbb C}\cong\bigoplus_{\pi\in\Pi_-}\pi_2^{U_2}\otimes\pi_5^{U_5}$ with $\Pi_-$ the representations with
$\pi_\infty$ trivial, trivial central character, $\pi_2^{U_2}\ne0$, $\pi_5^{U_5}\ne0$, spherical away from $2,5,13$ and
$\pi_{13} = \mathrm{St}$ (reversal $-1$); $T|_\pi = a_{17}(\pi)$, $m(\pi) = \dim\pi_2^{U_2}\dim\pi_5^{U_5}$, $\sum m(\pi) = 721$; $\Pi_+$ has
the quadratic twist. By Jacquet--Langlands these are weight-two cuspidal eigensystems with discrete series at $2$ and
$\mathrm{St}$ at $13$, and $|a_{17}|\le2\sqrt{17}$ is Eichler--Shimura--Igusa plus Weil for curves (no Deligne); all external
theorems not byte-cited. Finite consistency: $\mathrm{PGL}_2(\mathbb F_5)\cong S_5$ has irreducible degrees $1,1,4,4,5,5,6$ (seven,
not eight), $K_{\mathbb C}\cong6\,\mathbb C[G]\oplus1$ as a $G$-module (isotypic multiplicities $7,6,24,24,30,30,36$), every
$T$-eigenspace a $G$-module; the integer eigenvalues $7,-7,-1,-3,4,-5,6,-6,2,3,-2$ have multiplicities
$4,4,8,6,18,12,15,17,20,44,77$ ($225$ dimensions; the printed factor $x+5$ is the eigenvalue $-5$), trace $18$. An integer
$a_{17}$ does not identify a rational newform or an elliptic curve; the individual eigensystems are open.
\end{theorem}

\subsection{Frobenius, pairings, and the metric cone}

\begin{theorem}[The Bass doubling is Frobenius at $17$ as a one-operator system; the pairing is additional data]
\label{thm:lps-bass-frobenius}
\claimstatus{thm:lps-bass-frobenius}{proved}
(a) With $V_\pi$ the two-dimensional weight-two Galois representation of $\pi\in\Pi_-$ (geometric Frobenius, polynomial
$X^2-a_{17}X+17$, roots distinct by the strict bound), $(K\oplus K,F_1)\cong\bigoplus_{\Pi_-}(V_\pi,\mathrm{Frob}_{17})^{\oplus m(\pi)}$
over an algebraically closed field, as one-operator semisimple systems ($\alpha+\beta = a_{17}$, $\alpha\beta = 17$, unitary Satake
parameters $\alpha/\sqrt{17}$): a Hecke-selected summand of $H^1$ of modular Jacobians, not the whole $H^1$ of a genus-$721$
curve; adjoining $1$ and $17$ gives a curve-shaped package, not the zeta of that curve. (b) Jacquet--Langlands identifies
representations and Hecke actions, not the doubled graph basis with integral cycles; Eichler--Shimura's period isomorphism
($S_2\oplus\bar S_2\cong H^1(C,\mathbb C)$, Petersson $= \tfrac i2\int\omega\wedge\bar\omega$) supplies a Betti Hodge structure but no
real operator $\mathrm{Frob}_{17}$ preserving it: the brief's ``$\Omega_B$ and $G$ are the cup and Petersson pairings'' is % bare-ok
refuted as an implication and \emph{open} as a comparison. (c) On the graph model $F_1^{t}\Omega_BF_1 = 17\Omega_B$, % bare-ok
$F_1^{t}GF_1 = 17G$, $G>0\iff\|T\|<2\sqrt{17}$ (equality is a Jordan block: no positive invariant metric), and on an
$a$-eigenspace with $r_a = \sqrt{17-a^2/4}$, $J_a = (F_a-a/2)/r_a$, $\Omega_BJ_a = G_a/r_a$: the certificate $G$ and the % bare-ok
compatible metric $\Omega_BJ = G/r(T)$ differ (the brief's $G = \Omega_BJ$ is refuted); for fixed $\Omega_B,F_1$ the compatible % bare-ok
complex structure is unique. Total $H^1(X) = 0$ and $\dim K$ odd exclude an alternating form from the faces or on $K$ alone:
doubling bypasses both obstructions without a geometric comparison.
\end{theorem}

\begin{proposition}[The invariant metric cone has dimension $16539$; no symmetric ray is supplied]
\label{prop:lps-metric-cone}
\claimstatus{prop:lps-metric-cone}{proved}
$F_1^{*}HF_1 = 17H$ reads $(\bar\lambda_i\lambda_j-17)H_{ij} = 0$ in an eigenbasis, so $H$ is block-diagonal with an arbitrary positive
Hermitian block on each repeated eigenvalue, the real structure relating the $\alpha$ and $\beta$ blocks by conjugation: in the
two-copy coordinates $H_a = \begin{psmallmatrix}A&-aA/2+N\\-aA/2-N&17A\end{psmallmatrix}$, $A^{t} = A$, $N^{t} = -N$, one copy of
$\mathrm{Herm}^+_{m_a}$ per distinct $a$, real dimension $\sum_a m_a^2 = 16539$ (the brief's ``diagonal, one scalar per copy'' is
refuted). In the raw companion normalisation $e_\alpha = (\alpha v,v)$ the weights are $2r_a^2$ for $G$ and $2r_a$ for $\Omega_BJ$, not % bare-ok
constant across $a$; the exit-normalised basis and the transitive symmetries of \cref{thm:arithmetic-metric} are not data of
this model, so its one-ray conclusion cannot be imported. If a Hecke-equivariant comparison to weight-two forms is fixed,
Petersson gives $G_P = \begin{psmallmatrix}P&-PT/2\\-PT/2&17P\end{psmallmatrix}$ with $P = p_f$ on a line, $p_f = \|f\|^2_{\mathrm{Pet}}$
for a unit graph vector mapped to the normalised $f$; the ratio $G/G_P$ depends on the comparison and is open
(\texttt{notes/cerednik-drinfeld/}).
\end{proposition}

\subsection{The Weil--LPS channel and CP realisation}

\begin{theorem}[The Weil--LPS $(13;17)$ channel is a different Fourier selection of the same Hecke family]
\label{thm:weil-lps-versus-square-complex}
\claimstatus{thm:weil-lps-versus-square-complex}{proved}
For $(p;q) = (13;17)$ the Cayley group of \cref{sec:weil-lps} is $G^+ = \mathrm{PSL}_2(\mathbb F_{13})$ (order $1092$) and
$\Sigma_\pm = 18\Phi_\pm = \sum_{s\in S}\mathrm{End}(W_\pm)(s)$; finite Fourier decomposition $\mathbb C[G^+] = \bigoplus_\sigma\sigma\otimes M_\sigma$ gives
the spectral containment and selects multiplicity spaces $M_\sigma$; adelically it is the identity-determinant component of
$\mathcal A(U_2K(13))$ (principal level $13$, maximal at $5$), with $\mathrm{mult}_{\Sigma}(\lambda) = \sum_\sigma n_\sigma\dim\mathrm{Hom}_{G^+}(\sigma,\mathcal A^+_\lambda)$.
$W_-$ ($\dim6$) is a finite cuspidal Weil constituent (G\'erardin 1977, not byte-cited), not Steinberg ($\dim13$), and
$\mathrm{End}(W_-) = 1\oplus\tau_7\oplus\tau_{14,1}\oplus\tau_{14,2}$ (character $36,0,1,(7\pm\sqrt{13})/2$ on the classes), which contains no
finite Steinberg type: the channel does not select the Steinberg-at-13 forms of \cref{thm:lps-spectrum-hecke}. Exactly,
$\chi_{\Sigma_-}(x) = (x-18)(x-3)^8(x+3)^3(x+5)^6(x^2-3x-2)^3(x^4-34x^2+97)^3$, whose gcd with $\chi_T$ is
$(x-3)^8(x+3)^3(x+5)^6$; neither divides the other, and Bass roots cannot repair this. Harrow containment is Fourier
decomposition, not gradient cancellation; the ungraded odd-Weil Hashimoto zeta has its $70$ nontrivial Bass roots as poles,
whereas the worked model puts its $1442$ in the numerator. The brief's identification is refuted; the two are selections
of one spherical Hecke family at different levels.
\end{theorem}

\begin{observation}[A graded CP realisation of the net divisor is open; the shortcuts fail]
\label{obs:lps-graded-cp-open}
\claimstatus{obs:lps-graded-cp-open}{open}
For a graded transfer channel (\cref{def:graded-transfer-channel}) $\str\Edb^n = \sum_w|\Tr(PA_w)|^2\ge0$, $\Tr\Edb^n\ge0$ and
$\dim\mathrm{End}(V)_{\mathrm{ev}}-\dim\mathrm{End}(V)_{\mathrm{odd}} = (\dim V_+-\dim V_-)^2\ge0$; the sector prescription $\Esec0 = 1\oplus17$,
$\Esec1 = F_1$ has $\Tr\Edb^2 = 290-11230 = -10940<0$ and is impossible. Equality of supertraces fixes only the net divisor
$[1]+[17]-\sum[\lambda]$ with $\nu(0) = (\dim V_+-\dim V_-)^2+1440$; $N_n\ge0$ for all $n$ ($N_n\ge1+17^n-1442\cdot17^{n/2}$ for
$n\ge6$), so no obstruction from negativity; a homogeneous Kraus family realising this divisor with cancelling modes is
neither constructed nor excluded. Scalar orientation signs cancel under $\mathrm{Ad}$ ($\mathrm{Ad}(-U) = \mathrm{Ad}(U)$), so the
signed non-backtracking operator is a spectral, not a CP, transfer: $N_n = \Tr B_v^n-\Tr B_e^n+5760(1+(-1)^n)$. A CP
transfer with a letter register, $A_{(j,i)} = U_i\otimes|j\rangle\langle i|$ ($j\ne\mathrm{rev}(i)$), realises the unitary Hashimoto
block map of \cref{sec:weil-lps} (and its graded version for homogeneous $U_i$) with zero register coherences; the
script's $M_6$ carries the trivial doubled parity, the full $\mathbb C^7\oplus\mathbb C^6$ bond has sectors $85,84$. Hilbert--Schmidt
self-adjointness does not supply the companion metric: $F_1^{*}F_1\ne17$, and $G$ is positive only after the strict band
(\cref{obs:kraus-dichotomy}).
\end{observation}

\begin{numerical}[Blind numerics lane, the LPS square complex]
\label{num:lps-square-complex}
\claimstatus{num:lps-square-complex}{numerical}
\texttt{scripts/lps\_square\_complex.py} ($50$ assertions, $58$ s, numpy/scipy, exact integer arithmetic where the claim is exact;
\texttt{outputs/lps\_square\_complex.txt}; ledger D01--D50 and findings F1--F5 in \texttt{notes/deninger-lps/numerics.md}; $48$ pass,
the two failures being corrections of the brief, not of the example): sizes, both adjacency spectra with multiplicities,
$252$ unique rules, $7560$ squares with four representatives each, $B_1B_2 = 0$, ranks $119$ and $1801$ (modulo $101$ plus
$B_1B_2 = 0$), Betti $(1,0,5759)$, $\dim K = 721$, $T_1$ symmetric, $T_1D = DA_{17}$, $\|T\| = 4+2\sqrt3$, the full factor table
(deviation $5\cdot10^{-14}$), $\chi_T = \chi_{T_1}(x-18)/\chi_{A_{17}}$, $F_1^{t}\Omega F_1 = 17\Omega$, $F_1^{t}GF_1 = 17G$, $G>0$ (least % bare-ok
eigenvalue $0.172$), $|\mathrm{eig}F_1| = \sqrt{17}$, $N_1..N_6 = 0,11520,10080,126720,1209600,23886720$; $307$ norm-$289$ elements
and $T_{17}^2-17 = T_{289}$ on vertices; on edges $T_1^2-17 = (\text{double transport over }306) + 1$; the reversal-symmetric
sector carries $-\mathrm{spec}\,T$ (F1); $\mathrm{PGL}_2(\mathbb F_5)$ has seven irreducibles (F2); per-eigenvalue representation content
(e.g. $-2$: $4\times3+4'\times5+5\times9$; $3$: $1\times4+4\times5+5\times4$). The review reproduced $\sum m_a^2 = 16539$,
$\Omega J = G\,\mathrm{diag}(R^{-1},R^{-1})$ to $10^{-8}$, $\Tr F_1^2 = -11230$, and $\chi_{\Sigma_-}$ exactly. % bare-ok
\end{numerical}

\begin{observation}[What this changes; next]
\label{obs:deninger-lps-next}
\claimstatus{obs:deninger-lps-next}{sketched}
The finite Deninger package's positivity is explained arithmetically (weight-two Ramanujan for the selected eigensystems)
but not derived from a surface: the cup/Petersson comparison is the missing datum, and the natural candidate is
Cerednik--Drinfeld (the $13$-graph as the dual graph of the special fibre of a Shimura curve, the pairing as Grothendieck's
monodromy pairing; lane \texttt{notes/cerednik-drinfeld/}). The two-prime orbit quotient is non-Hausdorff
(\cref{obs:gl1-two-prime-solenoid}); the notebook's Weil--LPS channel is a sibling selection, not the same object; an exact
graded CP realisation of the net divisor is the open inverse problem.
\end{observation}
